\documentclass[11pt,a4paper]{article}

% ──────────────────────────────────────────────────────────────────────────
% EasyNet Ontology & CLI Layering — Consensus Specification
% Status: Draft for review before implementation
% Author: Silan Hu
% Build:   xelatex easynet_ontology.tex   (CJK + Unicode required)
% ──────────────────────────────────────────────────────────────────────────

\usepackage[margin=1in]{geometry}
\usepackage{amsmath,amssymb,amsthm}
\usepackage{xcolor}
\usepackage[colorlinks=true,linkcolor=blue!60!black,citecolor=blue!60!black,urlcolor=blue!60!black]{hyperref}
\usepackage{booktabs}
\usepackage{longtable}
\usepackage{array}
\usepackage{tabularx}
\usepackage{xltabular}
\newcolumntype{L}[1]{>{\raggedright\arraybackslash}p{#1}}
\usepackage{listings}
\usepackage{enumitem}
\usepackage{tcolorbox}
\usepackage{titlesec}
\usepackage{fancyhdr}
\usepackage{parskip}
\usepackage{xspace}

% ── Unicode + CJK support (xelatex) ───────────────────────────────────────
\usepackage{fontspec}
\usepackage{xeCJK}
% Pick a CJK font that ships with macOS by default.
\setCJKmainfont{PingFang SC}[
    BoldFont={PingFang SC Semibold},
    ItalicFont={PingFang SC Light}
]

% ── Page style ────────────────────────────────────────────────────────────
\pagestyle{fancy}
\fancyhf{}
\fancyhead[L]{\small EasyNet Ontology \& CLI Layering}
\fancyhead[R]{\small Consensus Specification}
\fancyfoot[C]{\thepage}
\renewcommand{\headrulewidth}{0.4pt}

% ── Section styling ───────────────────────────────────────────────────────
\titleformat{\section}{\Large\bfseries\color{black!85}}{\thesection}{0.7em}{}
\titleformat{\subsection}{\large\bfseries\color{black!75}}{\thesubsection}{0.6em}{}
\titleformat{\subsubsection}{\normalsize\bfseries\color{black!70}}{\thesubsubsection}{0.5em}{}

% ── Tag colors ────────────────────────────────────────────────────────────
\definecolor{axiomcolor}{RGB}{180,30,30}
\definecolor{derivedcolor}{RGB}{30,90,180}
\definecolor{deferredcolor}{RGB}{120,90,30}
\definecolor{codebg}{RGB}{245,245,248}
\definecolor{codeborder}{RGB}{210,210,220}

% ── Tag macros ────────────────────────────────────────────────────────────
\newcommand{\axiom}{\textcolor{axiomcolor}{\textbf{[axiom]}}\xspace}
\newcommand{\derived}{\textcolor{derivedcolor}{\textbf{[derived]}}\xspace}
\newcommand{\deferred}{\textcolor{deferredcolor}{\textbf{[deferred]}}\xspace}

% ── Code listing style ────────────────────────────────────────────────────
\lstdefinestyle{eal}{
    backgroundcolor=\color{codebg},
    basicstyle=\ttfamily\small,
    frame=single,
    rulecolor=\color{codeborder},
    breaklines=true,
    showstringspaces=false,
    keywordstyle=\color{blue!70!black}\bfseries,
    commentstyle=\color{green!40!black}\itshape,
    stringstyle=\color{red!60!black},
    morekeywords={class,public,private,ability,skill,let,print,return,use},
    columns=fullflexible,
    keepspaces=true,
    aboveskip=0.6em,
    belowskip=0.6em,
}
\lstset{style=eal}

% ── Theorem-like environments ─────────────────────────────────────────────
\newtcolorbox{invariant}[1][]{
    colback=red!4!white,
    colframe=red!50!black,
    fonttitle=\bfseries,
    title={Invariant},
    sharp corners=south,
    boxrule=0.6pt,
    #1
}
\newtcolorbox{notebox}[1][]{
    colback=blue!4!white,
    colframe=blue!50!black,
    fonttitle=\bfseries,
    title={Note},
    sharp corners=south,
    boxrule=0.6pt,
    #1
}
\newtcolorbox{warningbox}[1][]{
    colback=orange!6!white,
    colframe=orange!70!black,
    fonttitle=\bfseries,
    title={Transitional misalignment},
    sharp corners=south,
    boxrule=0.6pt,
    #1
}

% ── Title ─────────────────────────────────────────────────────────────────
\title{
    \vspace{-2em}
    \textbf{\Large EasyNet Ontology and CLI Layering}\\[0.3em]
    \large Consensus Specification (pre-implementation draft)
}
\author{
    Silan Hu \\ \small National University of Singapore
}
\date{Draft, \today}

% ──────────────────────────────────────────────────────────────────────────
\begin{document}
\maketitle

\begin{abstract}
\noindent
This document records the consensus reached between the EasyNet authors and a
rapporteur during a multi-round design alignment session, regarding the
ontology that underpins EasyNet---an agent-native distributed system---and
the layering of its command-line interface. The current protocol-facing claim
is that an \emph{ability} is a network-visible executable contract with a
stable public surface, policy, Invocation, observability, and receipt
semantics. A \emph{skill} is implementation/resource-plane support: an
ability implementation may use skills, memory, workflow graphs, tools, or
other abilities, but a skill is not itself network-addressable. We further
pin down the object-oriented view of agents (public abilities, private
skills), the two-language two-runtime split (AAL/BCC for production;
EAL/EasyNet for consumption), the demotion of \emph{device} from a
network-first-class object to a hosting substrate, and the desugaring of
\texttt{agent send} into a single-line External EAL mission invoking the
target's default \texttt{chat} ability. Each statement in the document is
explicitly tagged as an axiom (directly stated by the principals), a
derivation (logically deduced from one or more axioms), or a deferred item
(known gap, deliberately postponed). The intent of this artifact is to lock
the model before any implementation work begins, so that the upcoming
refactor and all future epics share the same vocabulary.
\end{abstract}

\tableofcontents
\vspace{1em}
\hrule
\vspace{1em}

% ──────────────────────────────────────────────────────────────────────────
\section{Reading guide}
\label{sec:reading-guide}

Every substantive statement in this document carries one of three tags:

\begin{itemize}[leftmargin=2.4em,itemsep=0.3em]
    \item[\axiom] Directly stated by Silan, Sean, or in the alignment session
    transcript. Not negotiable in this document. If you wish to revise an axiom,
    a separate alignment round is required.

    \item[\derived] Logically derived from one or more axioms. Negotiable: if
    you disagree with a derivation, the disagreement should be raised
    \emph{before} implementation begins, and either the derivation is revised
    or the axioms it depends on are clarified.

    \item[\deferred] A known gap, deliberately left for a future epic. The
    document explicitly identifies these so that downstream design work does
    not silently fill them with the wrong assumption.
\end{itemize}

The tagging discipline is the most important property of this document. It
makes the difference between ``what we agreed'' and ``what I inferred''
auditable.

% ──────────────────────────────────────────────────────────────────────────
\section{The two core technical theses}
\label{sec:two-theses}

\subsection{Thesis A: Agent Ability}

\axiom{}
\begin{quote}
\textbf{Ability is the network-visible executable contract. Skill is private
implementation support.}

\medskip
\noindent
An ability exposes a stable public surface, policy boundary, Invocation
entrypoint, observability stream, and receipt semantics. Its implementation
may evolve through long-context memory, workflow graphs, model/tool use, skill
resources, and child Invocations, but those implementation resources are not
part of the public addressable contract.
\end{quote}

The older shorthand
``\emph{skill plus self-evolution plus network-wide invocation}'' was useful
as a draft intuition, but it is not the normative definition. It must not be
read as ``Skill is a component of Ability'' or ``private Ability equals
Skill.'' The protocol boundary is:

\begin{itemize}[leftmargin=2em,itemsep=0.3em]
    \item \textbf{Ability} --- public, addressable, policy-governed, invoked
    through Axon Invocation, and producing receipts.
    \item \textbf{Skill} --- private implementation/resource-plane knowledge
    that may support an ability but cannot be invoked across the network.
    \item \textbf{Implementation graph} --- memory, workflow, tools, skills,
    model calls, and child Invocations behind the stable ability surface.
\end{itemize}

\derived{} Removing one of the public-contract properties degenerates the
object into something strictly simpler:
\begin{itemize}[leftmargin=2em,itemsep=0.2em]
    \item without Invocation/receipt semantics $\Rightarrow$ local tool,
    \item without stable schema/policy $\Rightarrow$ ad-hoc handler,
    \item without observability/audit $\Rightarrow$ unauditable side effect.
\end{itemize}

\subsection{AbilityDescriptor, AuthorityBinding, and AbilityImpl}

\axiom{}
\begin{quote}
\textbf{AbilityDescriptor is the governed interface; AbilityImpl is the
executable binding; Receipt is the versioned, verifiable execution fact.}
\end{quote}

The object-oriented shorthand ``agent owns ability'' is not precise enough for
the device/daemon/plugin control plane. The cleaner split is:

\begin{itemize}[leftmargin=2em,itemsep=0.3em]
    \item \textbf{User Principal} --- accountability and authentication
    subject for a human or service account. A User account is not an Agent; it
    may own or authorize hosted Agents, but it does not advertise abilities as a
    synthetic account Agent.
    \item \textbf{Agent} --- routable behavior subject that may advertise
    governed \texttt{AbilityDescriptor}s and appear as Invocation callee. An
    Agent has logical identity independent of the Device or process that hosts
    its implementation.
    \item \textbf{Device} --- execution substrate identified by
    \texttt{device\_ura}, with \texttt{node\_id} and local resources.
    \item \textbf{Daemon} --- projection and dispatch runtime. It registers
    descriptors, binds implementations, dispatches Invocations, and emits
    receipts.
    \item \textbf{System Agent} --- restricted Agent sponsored by one Device
    for device-native abilities. It advertises the descriptor and remains the
    logical callee; the Device is its host and possible receipt attestor.
    \item \textbf{Authority} --- routable realm governance actor. It may own
    authority-scoped descriptors and issue or verify governance predicates, but
    it is not an execution substrate.
    \item \textbf{AbilityDescriptor} --- versioned governed interface with
    \texttt{ability\_ura}, \texttt{name}, \texttt{version}, \texttt{schema},
    \texttt{schema\_hash}, \texttt{call\_mode}, visibility, and policy
    reference.
    \item \textbf{AuthorityBinding} --- governance predicate over both
    advertisement and invocation. A binding that only proves publication
    rights is incomplete.
    \item \textbf{AbilityImpl} --- executable binding for a descriptor version,
    with implementation kind, entrypoint, \texttt{impl\_hash}, and
    \texttt{runtime\_env}.
    \item \textbf{PluginAbilityImpl} --- local plugin-provided execution
    binding loaded and bound by the daemon.
\end{itemize}

\derived{} Ownership, authority, and execution binding are separate layers.
For device-native behavior, the System Agent advertises and owns the
AbilityDescriptor. The Device is sponsor, host, key custodian, and possible
receipt attestor; it is not a Principal. Accountability resolves through the
Device's paired Principal and the explicit \texttt{AuthorityBinding}.

\begin{quote}
\textbf{Canonical actor contract.}
User/Account = Principal/accountability;
Agent = routable AbilityDescriptor owner/callee;
System Agent = Device-sponsored restricted Agent;
Authority = realm governance actor;
Device = execution host/custody/sponsor, not an ordinary public actor;
AbilityImpl/Plugin/Skill = implementation/resource plane, not public ownership identity.
\end{quote}

\derived{} \texttt{Invocation.subject} is an \texttt{EntityRef}, not only a
resource pointer. It may refer to a resource, agent, ability, session,
continuation, state object, or the callee itself.

\derived{} Versioning is part of the security model. An Invocation must pin
\texttt{descriptor\_version}. A Receipt must bind
\texttt{descriptor\_version}, \texttt{schema\_hash}, \texttt{impl\_hash},
\texttt{runtime\_env}, \texttt{authority\_proof}, \texttt{input\_hash},
\texttt{output\_hash}, \texttt{parent\_receipts}, and a signature. Otherwise
the system cannot prove which governed interface and executable binding were
used at execution time.

\begin{invariant}
Descriptor is not implementation. Advertisement is not authority. Authority is
not execution. Invocation is not a raw handler call. Receipt is not a return
value.
\end{invariant}

\subsection{Thesis B: EasyNet, the network of vocational abilities}

\axiom{}
\begin{quote}
\textbf{EasyNet is the communication-and-collaboration network for
network-wide vocational abilities.}
\end{quote}

Two interaction modalities live on EasyNet:

\begin{itemize}[leftmargin=2em,itemsep=0.3em]
    \item \textbf{Communication} (light interaction): post, feed, chat,
    connection.
    \item \textbf{Collaboration} (deep interaction): functional execution, the
    actual cross-agent ability calls.
\end{itemize}

\derived{} The CLI under design here is concerned almost exclusively with the
\emph{collaboration} modality. The communication modality is a separate
surface (most likely UI / social), and is intentionally out of scope for this
document.

% ──────────────────────────────────────────────────────────────────────────
\section{Self-evolution as graph}
\label{sec:graph}

\subsection{Internal structure of an evolving ability}

\axiom{} The core of self-evolution is a \emph{graph}, which is split into
\emph{memory} and \emph{workflow}:

\begin{center}
\begin{tabular}{l}
\texttt{Ability implementation = graph} \\
\quad\texttt{|-- memory graph}\quad (knowledge / long context / accumulated state) \\
\quad\texttt{|-- workflow graph}\quad (execution structure / decision flow / call topology) \\
\end{tabular}
\end{center}

\subsection{What ``self-evolution'' actually means}

\derived{}
\begin{itemize}[leftmargin=2em,itemsep=0.3em]
    \item The \emph{public surface} of an ability (its signature) does not
    evolve freely; that would break callers.
    \item The \emph{graph} behind the surface evolves continuously as new
    memory accrues (from long context, traces, and user feedback) and as new
    workflow paths are discovered (better orderings, new sub-ability calls).
    \item An ability is therefore \emph{the same ability} even as its internal
    graph grows. Distillation operates on the \emph{graph}, not on the surface
    signature.
\end{itemize}

\subsection{What distillation operates on}

\derived{}
\begin{quote}
Distillation only ever distills \textbf{abilities}, never skills.
\end{quote}

The reason is observability: skills are private to a single agent and have no
network trace; abilities are public and have invocation traces from across
the network. Trace data feeds back into the graph (memory growth and workflow
refinement). It does not feed back into skills.

% ──────────────────────────────────────────────────────────────────────────
\section{The OOP view (load-bearing)}
\label{sec:oop}

\subsection{Agent as object, ability as public method, skill as private field}

\axiom{} An agent class in AAL takes the following shape:

\begin{lstlisting}
class Agent {
    public:
        ability chat(context) -> reply
        ability review(code) -> review
        ability decompose(goal) -> [subgoals]

    private:
        skill thinking_style
        skill code_review_prior
        skill memory_compression_strategy
        skill ...
}
\end{lstlisting}

\axiom{} \texttt{chat} and \texttt{review} are typical \emph{default abilities}
every agent exposes, analogous to \texttt{Object.toString()} in Java.

\subsection{Visibility rules}

\axiom{}
\begin{itemize}[leftmargin=2em,itemsep=0.3em]
    \item \textbf{Ability is public.} It has a frozen signature, is
    discoverable network-wide, and is callable from External EAL.
    \item \textbf{Skill is private.} It has no signature on the network, is
    not discoverable from outside the agent, and is \emph{never} callable
    from External EAL.
    \item \textbf{EAL (external) can only write \texttt{agent.ability(...)}}.
    It cannot reach into skills.
    \item \textbf{Public calls go over gRPC.} That is what makes
    ``network-wide invocation'' real.
\end{itemize}

\subsection{Five questions collapsed into one axiom}
\label{sec:five-into-one}

The visibility rule simultaneously answers five design questions that, before
the alignment session, appeared independent:

\begin{center}
\renewcommand{\arraystretch}{1.25}
\begin{tabularx}{\linewidth}{L{0.45\linewidth} X}
\toprule
\textbf{Question} & \textbf{Answer (because of visibility rule)} \\
\midrule
Why does ability need a graph? &
Because it must self-evolve; static structure cannot grow. \\
Why is skill private? &
Because it is single-agent internal knowledge; sharing makes no sense across
agent boundaries. \\
Why gRPC? &
Because ability is a network service. \\
Why does the EAL compiler only see ability? &
Because it operates in \texttt{public} scope. \\
Why does distillation only distill abilities? &
Because only public calls produce observable network traces. \\
\bottomrule
\end{tabularx}
\end{center}

\derived{} These five answers are now consequences of one rule, not five
independent design decisions. This collapse is the central simplification of
the alignment session.

\subsection{Encapsulation invariant}

\begin{invariant}
No CLI command, no SDK call, and no EAL construct may reach across an agent
boundary into a skill. If a future need arises (``I want my agent to learn
from another agent's skill''), the only valid answer is: \emph{the other agent
must wrap that skill as an ability.}
\end{invariant}

\derived{} This invariant is the load-bearing wall of the system. Violating it
collapses the OOP model and, with it, the ability/skill distinction. Every
proposed CLI verb in this document is required to pass an encapsulation
check before being added.

% ──────────────────────────────────────────────────────────────────────────
\section{Two languages, two runtimes}
\label{sec:two-languages}

\subsection{The pairing}

\axiom{} EasyNet has two distinct language/runtime pairs:

\begin{center}
\renewcommand{\arraystretch}{1.3}
\begin{tabularx}{\linewidth}{L{0.27\linewidth} L{0.22\linewidth} X}
\toprule
\textbf{Language} & \textbf{Interpreter} & \textbf{Expresses} \\
\midrule
AAL (Agent Abstraction Language) & BCC (Brain Cell Coding) & how an agent
class is defined \\
EAL (EasyNet Ability Language) & EasyNet & how abilities compose to
accomplish a task \\
\bottomrule
\end{tabularx}
\end{center}

BCC interprets AAL. EasyNet interprets EAL. \textbf{They are two independent
language/runtime pairs}, not a single compilation pipeline.

\subsection{Producer side versus consumer side}

\derived{}
\begin{itemize}[leftmargin=2em,itemsep=0.3em]
    \item AAL/BCC is the \emph{producer side}: who builds an agent, what
    abilities it exposes.
    \item EAL/EasyNet is the \emph{consumer side}: who orchestrates agents to
    accomplish a task.
    \item The bridge between them is the \emph{gRPC ability endpoint}: BCC
    exposes abilities as gRPC services; EasyNet calls them via gRPC.
\end{itemize}

\derived{} This separation means: \textbf{writing a new agent does not
require knowing EAL; writing a new mission does not require knowing AAL.}
This is clean separation of concerns and is one of the strongest properties
of the design.

\subsection{Two layers of EAL}
\label{sec:two-eal-layers}

EAL appears in two distinct scopes:

\begin{center}
\renewcommand{\arraystretch}{1.25}
\begin{tabularx}{\linewidth}{L{0.16\linewidth} X X}
\toprule
 & \textbf{External EAL} & \textbf{Internal EAL} \\
\midrule
Who writes & Humans, in \texttt{*.eal} files & Generated by the ability graph
at call time \\
Scope & \texttt{public} --- only sees abilities & \texttt{private} --- sees
the agent's own skills, memory, abilities \\
Enforced by & EasyNet's EAL parser at compile time & The agent's own runtime \\
Persisted as & Mission runs & Per-call execution traces inside an ability
invocation \\
OOP analogue & Application code calling public methods & The body of a
method, which can touch private fields \\
\bottomrule
\end{tabularx}
\end{center}

\derived{} The two layers use the same surface language but operate under
different access scopes, exactly as Java syntax compiles differently inside
versus outside a class definition.

\subsection{Where the current \texttt{easynet} binary fits}

\derived{} The current \texttt{easynet} binary is the \emph{local CLI surface
for the EAL runtime} (i.e.\ EasyNet itself). It does not yet contain BCC or
any AAL machinery. Anywhere this document refers to ``the EasyNet runtime,''
the local artifact in question is the \texttt{easynet} binary plus its
embedded \texttt{axon} runtime process.

\deferred{} BCC will arrive as a separate concern. When it does, the
\texttt{easynet} binary may grow a sibling (\texttt{bcc}) or absorb BCC as a
subcommand. That decision is future work.

% ──────────────────────────────────────────────────────────────────────────
\section{Device is a hosting substrate (interpretation C)}
\label{sec:device-c}

\subsection{Device is not a network-first-class entity}

\axiom{}
\begin{quote}
Device is \textbf{not} a first-class network object. It is the physical
substrate on which abilities are deployed --- like an OS to a process.
\end{quote}

Three interpretations of ``device'' were considered during the alignment
session:

\begin{itemize}[leftmargin=2em,itemsep=0.3em]
    \item[(A)] Device $=$ degenerate Agent (an agent without any skill).
    \emph{Rejected}: pulls device back into network identity.
    \item[(B)] Device $=$ a parallel specialization of Process, peer to
    Agent. \emph{Rejected}: same problem.
    \item[(C)] Device $=$ hosting substrate, not first-class on the network.
    \textbf{Chosen.}
\end{itemize}

\subsection{What this means concretely}

\derived{} The network-first-class actor set is deliberately small:

\begin{center}
\renewcommand{\arraystretch}{1.25}
\begin{tabularx}{\linewidth}{X X}
\toprule
\textbf{Network first-class actor/interface} & \textbf{Host-local, not first-class actor} \\
\midrule
Agent, System Agent, and Service (routable callees) & Device (physical
substrate, hardware, credentials) \\
Authority (realm governance actor) & Axon runtime (the protocol/runtime
process) \\
AbilityDescriptor (governed callable interface) & MCP server (local stdio process for AI
clients) \\
\bottomrule
\end{tabularx}
\end{center}

\derived{} Ordinary public Invocation callees are \textbf{Agent, System Agent,
Service, or Authority} identities advertising governed AbilityDescriptors.
Devices exist as execution host, sponsor, key custodian, and possible receipt
attestor. A Device may appear in bounded bootstrap, pairing,
publication-custody, or session-control caller paths enumerated by the shared
\texttt{DeviceCallerPurpose} classifier, but it is not an ordinary public
callee.

\subsection{Logical identity versus physical placement}

\derived{}
\begin{itemize}[leftmargin=2em,itemsep=0.4em]
    \item \textbf{Logical identity} (network-visible): an agent is identified
    by \texttt{<tenant>/<agent-name>}, e.g.\ \texttt{silan/reviewer}. This is
    what appears in \texttt{agent.ability(...)} calls.
    \item \textbf{Physical placement} (hidden): which device currently hosts
    that agent is metadata, accessible only through \texttt{device} CLI
    commands and only for diagnostic purposes.
\end{itemize}

The current PR uses \texttt{<tenant>/<agent-name>} as the minimum viable form.
A future extension to \texttt{<tenant>/<agent-name>\#<instance-id>} is planned
for when multiple instances of the same agent class need to coexist.

\paragraph{Implementation Note.}
The current implementation of the agent identity layer is specified in
\texttt{docs/AGENT\_IDENTITY.md}. This document is a derived engineering
artifact (L2 identity layer) and is not part of the ontology consensus.
It may evolve independently as long as it does not contradict the logical
identity model described in this section and the deferred tenant semantics
in \S\ref{sec:gaps}.

\subsection{Cardinality}

\derived{} Because device is not first-class:
\begin{itemize}[leftmargin=2em,itemsep=0.3em]
    \item One device can host multiple agent instances.
    \item One agent instance can in principle span multiple devices (if its
    memory graph and workflow runtime live on different machines), and the
    network still sees a single agent.
\end{itemize}

\derived{} This is the affordance interpretation C buys us that A and B do
not. It \emph{leaves room for distributed agents}, even though the current
implementation does not yet exercise it.

\subsection{Known transitional surface}

\begin{warningbox}
The CLI surface \texttt{easynet ability deploy <path> --to <node>} still uses a
Device selector to choose the execution host. This is a host-selection UX, not
descriptor ownership. The current daemon implementation derives deployed
AbilityDescriptor ownership from the device-sponsored ability-management
System Agent while keeping the target Device as execution host.

\smallskip
\deferred{} A future UX may replace \texttt{--to <node>} with an Agent or
SystemAgent placement selector. Until then, code and documentation must keep the
three facts separate: selected host Device, descriptor owner SystemAgent, and
accountability User Principal.
\end{warningbox}

% ──────────────────────────────────────────────────────────────────────────
\section{Default abilities and call sugar}
\label{sec:default-call}

\subsection{\texttt{chat} as the default callable}

\axiom{} Every agent class typically exposes a \texttt{chat(prompt) -> reply}
method as a ``default callable.'' This is analogous to:

\begin{itemize}[leftmargin=2em,itemsep=0.2em]
    \item \texttt{Object.toString()} in Java,
    \item \texttt{\_\_call\_\_} in Python,
    \item \texttt{operator()} in C\texttt{++},
    \item the \texttt{()} of a Functor in Haskell.
\end{itemize}

\subsection{\texttt{agent send} is sugar for the default ability}

\axiom{}
\begin{quote}
``实际上 \texttt{send} 是默认交给 \texttt{agent.chat()} ability,只是默认语法糖。''
(``In fact, \texttt{send} dispatches to the \texttt{agent.chat()} ability by
default; it is merely default syntactic sugar.'')
\end{quote}

Two layers of sugar are at play:

\begin{lstlisting}
# Layer 1 (CLI sugar):
easynet agent send claude "hello"
  desugars to a single-line External EAL mission:
let r = claude.chat(prompt: "hello")
print(r)

# Layer 2 (EAL sugar, future):
let r = claude.chat(prompt: "hello")
  may be writable as:
let r = claude("hello")
  because chat is the agent's default callable
\end{lstlisting}

\derived{} Three consequences follow:

\begin{enumerate}[leftmargin=2em,itemsep=0.3em]
    \item \textbf{\texttt{agent send} is not a special command.} It is the
    shortest possible mission. This unifies ``send to one agent'' and ``run a
    mission'' under one execution model.
    \item \textbf{Agents are callable objects.} The default \texttt{chat}
    ability gives every agent a built-in call operator. Custom agents can
    override or supplement it.
    \item \textbf{Mission is the default orchestration implementation model.}
    Once \texttt{agent send} is sugar for a one-line mission, the CLI, EAL, and
    MCP-facing paths share the mission runtime for composition. Mission still
    emits ordinary child Invocations; it does not replace the Invocation
    primitive or redefine caller/callee/subject/nonce/causal-context semantics.
\end{enumerate}

% ──────────────────────────────────────────────────────────────────────────
\section{Six locked-in design decisions}
\label{sec:six-decisions}

These were resolved during the alignment session and are not reopened by this
document. They represent the minimum set of decisions required to begin
implementation.

\renewcommand{\arraystretch}{1.35}
\begin{xltabular}{\linewidth}{c L{0.24\linewidth} X X}
\toprule
\textbf{\#} & \textbf{Question} & \textbf{Decision} & \textbf{Rationale} \\
\midrule
\endfirsthead
\toprule
\textbf{\#} & \textbf{Question} & \textbf{Decision} & \textbf{Rationale} \\
\midrule
\endhead

1 & Agent instance ID format &
\texttt{<tenant>/<agent-name>}, with \texttt{\#<instance-id>} reserved for the
future &
Logical, network-visible, decoupled from device. \\

2 & Agent instantiation in current PR &
Not done. Agents remain implicit, provisioned by runtime plus deploy. &
AAL/BCC do not exist yet; premature commands would lock in wrong semantics. \\

3 & \texttt{device join} semantics &
Registers a hosting substrate (physical device plus credentials). &
Matches the Docker-daemon-register / Kubernetes-node-join pattern under
interpretation C. \\

4 & \texttt{agent send} semantics &
Sugar for a single-line External EAL mission invoking the target's default
\texttt{chat} ability. &
Unifies all cross-agent interaction under the mission execution model. \\

5 & \texttt{discuss} / \texttt{think} placement &
Primary location is \texttt{mission} (\texttt{mission discuss},
\texttt{mission think}). Deprecated alias retained under \texttt{agent} for
backwards compatibility. &
Both are mission patterns (multi-agent orchestration; iterative planning
loop), not agent-instance methods. \\

6 & \texttt{device config} versus \texttt{runtime config} &
Semantically \texttt{runtime config}; physically remains \texttt{device
config} for now. Long-help explicitly notes the migration. &
Settings (exec enable, MCP behavior) control the runtime, not the physical
device. Code migration deferred to avoid CLI churn. \\

\bottomrule
\end{xltabular}

% ──────────────────────────────────────────────────────────────────────────
\section{CLI surface (consensus shape)}
\label{sec:cli-surface}

\subsection{Top-level groups}

\begin{lstlisting}
easynet
|
|-- agent       Manage agent instances (network actors)
|     |-- add / list / remove / doctor       (instance lifecycle)
|     |-- send                                (sugar for chat)
|     |-- session new/list/show/append/end    (memory dimension)
|     |-- discuss          DEPRECATED -> use `mission discuss`
|     `-- think            DEPRECATED -> use `mission think`
|
|-- device      Manage hosting substrates (physical machines)
|     |-- join / reset / config               (this host's lifecycle)
|     `-- list / show / remove                (substrates on the network)
|         (no rename / tag: would require an admin ability that doesn't exist)
|
|-- ability     Manage public ability endpoints
|     |-- list / show                          (discovery)
|     |-- deploy / uninstall                   (lifecycle; deploy is host-routed)
|     |-- invoke                               (direct call)
|     `-- exec                                 (one-shot ad-hoc command)
|         (no `update`: conflates three time scales)
|         (no `logs`: wrong artefact)
|
|-- mission     Run External EAL programs and patterns
|     |-- compile / run / list / show / cancel
|     |-- discuss                              (multi-agent orchestration)
|     `-- think                                (iterative planning loop)
|
|-- runtime     Local Axon process lifecycle
|     `-- start / stop / status / connect / logs
|
|-- mcp         Local MCP server process
|     `-- serve / status
|
|-- doctor      Aggregated health check                          (NEW top-level)
`-- completion  Shell completion generator                       (NEW top-level)
\end{lstlisting}

\subsection{Backwards compatibility}

All current top-level commands (\texttt{start}, \texttt{stop}, \texttt{status},
\texttt{connect}, \texttt{devices}, \texttt{abilities}, \texttt{deploy},
\texttt{invoke}, \texttt{exec}, \texttt{mission}, \texttt{join}, \texttt{reset},
\texttt{config}, \texttt{mcp-server}, \texttt{mcp-install}, \texttt{agent},
\texttt{discuss}, \texttt{skill-install}, \texttt{think}) remain available as
\textbf{deprecated aliases} that:

\begin{itemize}[leftmargin=2em,itemsep=0.2em]
    \item continue to function identically,
    \item emit a one-line stderr deprecation notice pointing at the new
    layered command,
    \item will be removed in a future release (no date set).
\end{itemize}

This makes the current PR a \textbf{purely additive, zero-breakage} change.

\subsection{Deliberately absent commands}

The following commands are missing \emph{on purpose}. Adding them would
violate the ontology:

\begin{itemize}[leftmargin=2em,itemsep=0.35em]
    \item \texttt{easynet skill <anything>} --- skills are private; exposing
    them breaks the encapsulation invariant of \S\ref{sec:oop}.

    \item \texttt{easynet capability <anything>} --- capability is the
    cross-process invocation abstraction; introducing it half-baked would
    lock in premature semantics.

    \item \texttt{easynet ability update} --- conflates three distinct time
    scales (signature bump, graph evolution, policy retraining). See
    \S\ref{sec:time-scales}.

    \item \texttt{easynet ability logs} --- the real artefact is the ability
    graph trace, which lives at a layer the current PR does not yet model.

    \item \texttt{easynet agent spawn} --- agent instantiation is an AAL/BCC
    concern; nothing in the current PR can honestly spawn an agent class
    instance.
\end{itemize}

% ──────────────────────────────────────────────────────────────────────────
\section{Non-CLI artifacts shipped in this PR}
\label{sec:non-cli}

\subsection{The \texttt{ability\_graph\_traces} field on mission run metadata}

A new optional field is added to \texttt{MissionRunMeta}:

\begin{lstlisting}
/// Per-cross-agent-ability-call execution summaries. Each entry
/// captures what the target agent's ability graph did to satisfy
/// one call. Empty for runs that only invoked device abilities
/// (which have no graph).
///
/// Schema is intentionally Value here: this is a landing slot for
/// the upcoming ability-graph trace format, not the format itself.
#[serde(default, skip_serializing_if = "Option::is_none")]
pub ability_graph_traces: Option<Vec<serde_json::Value>>,
\end{lstlisting}

\derived{} Section~\ref{sec:graph} states that ability implementations are
graphs. The mission run metadata is the natural landing spot for traces of
those graphs, even before the trace format is fully specified. Naming the
field \texttt{ability\_graph\_traces} (instead of, e.g.,
\texttt{internal\_eal\_summaries} or \texttt{dispatch\_traces}) \emph{teaches
the next reader} that an ability has a graph, by the field name alone.

\subsection{Retention of \texttt{agent\_sessions.rs}}

The previously-considered deletion of this file is reversed. Sessions are the
simplest form of the \emph{memory} dimension of an agent (per-caller
conversation history). They live on the agent side of the OOP wall, are
private by default, and are properly modeled as ``the calling client's view
of one slice of an agent's memory.''

% ──────────────────────────────────────────────────────────────────────────
\section{Three time scales}
\label{sec:time-scales}

When evaluating any future CLI command, ask which time scale it operates on:

\begin{center}
\renewcommand{\arraystretch}{1.3}
\begin{tabularx}{\linewidth}{L{0.20\linewidth} L{0.20\linewidth} X L{0.20\linewidth}}
\toprule
\textbf{Time scale} & \textbf{Frequency} & \textbf{Example operation} &
\textbf{Allowed verb} \\
\midrule
Schema/SLA bump & discrete, human-in-loop & publishing v2 of an ability with a
new input field & \texttt{ability deploy} (new version) \\
Graph evolution & low-frequency, semi-automatic & a new memory pattern
accrues; a workflow path is shortened & none yet (internal) \\
Per-call execution & realtime & one ability call picks a route through the
graph and executes & \texttt{ability invoke}, \texttt{mission run} \\
\bottomrule
\end{tabularx}
\end{center}

\derived{} Verbs that try to compress two scales into one (for example
\texttt{ability update} covering both v-bump and graph tweak) \textbf{must be
rejected}. They look convenient but corrupt the user's mental model.

% ──────────────────────────────────────────────────────────────────────────
\section{Higher-level judgment and next deliverables}
\label{sec:next}

\axiom{} (Silan): the next deliverable beyond this document should be a
paper-level \textbf{System Model: Definition + Invariants + Diagram},
suitable for a whitepaper or top-tier venue framing. This is the system's
true moat.

The current PR explicitly \emph{does not} attempt to be that paper. It
implements the operational consequences of the consensus reached above. The
paper-level model is tracked as the next epic. A short pointer to that
deliverable will appear at the end of \texttt{docs/ARCHITECTURE.md}:

\begin{notebox}
\textbf{Next:} \texttt{docs/SYSTEM\_MODEL.md} (paper-level formalization).
Until that exists, this document and \texttt{ARCHITECTURE.md} are the source
of truth for ontological questions about EasyNet.
\end{notebox}

% ──────────────────────────────────────────────────────────────────────────
\section{Open recursion points}
\label{sec:gaps}

These are honestly identified gaps. They are not blockers for the current PR,
but they are blockers for the next epic.

\begin{enumerate}[leftmargin=2em,itemsep=0.35em]
    \item \textbf{Capability abstraction.} ``Capability'' as the unifying
    invocation interface across agent and device is mentioned but not
    specified. This is the next ontological extension.

    \item \textbf{Internal EAL trace format.} \S\ref{sec:two-eal-layers}
    states that internal EAL is persisted as per-call execution traces, but
    the schema is not specified. The \texttt{ability\_graph\_traces} field
    lands the slot but defers the format.

    \item \textbf{Default ability set.} \S\ref{sec:default-call} says
    \texttt{chat} is a typical default. Whether \texttt{chat} is \emph{the}
    default, whether \texttt{review} is also default, and whether agents
    must declare a default --- all unspecified.

    \item \textbf{Agent class system.} AAL is named, but its grammar, type
    system, and inheritance model (if any) are entirely unspecified.

    \item \textbf{Tenant model.} \texttt{<tenant>/<agent-name>} is chosen as
    the ID format but tenant lifecycle, sharing, and authorization are
    unspecified.

    \item \textbf{Status of \texttt{mcp-install} and \texttt{skill-install}.}
    Whether these legacy top-level commands should remain, move under
    \texttt{agent install}, or be removed entirely depends on whether
    skill-install should exist at all under interpretation C. They are kept
    as legacy aliases in the current PR with the question explicitly
    flagged.
\end{enumerate}

% ──────────────────────────────────────────────────────────────────────────
\section{Conclusion}

This document closes a multi-round alignment session by recording, with
auditable provenance, what the principals agreed and what the rapporteur
inferred. The single most consequential agreement is the OOP visibility
rule of \S\ref{sec:oop}, which by itself collapses five previously-open
design questions (\S\ref{sec:five-into-one}) into a single axiom. Around it,
the two-language two-runtime split (\S\ref{sec:two-languages}) and the
demotion of device to a hosting substrate (\S\ref{sec:device-c}) provide the
remaining structure required to begin implementation without anticipating
the next ontological extension.

The implementation that follows this document is, by design, the smallest
change that can reflect the new ontology in the CLI surface. Anything that
would require the next epic --- capability abstraction, AAL/BCC, the formal
system model --- is deliberately deferred and explicitly enumerated in
\S\ref{sec:gaps}, so that no future contributor will be forced to guess
which concepts are settled and which are still open.

\vfill
\hrule
\vspace{0.4em}
\noindent\small
\textit{This document is the consensus output of the alignment session and
should be reviewed in full before any code is written. If you disagree with
a \textbf{[derived]} statement, raise it before implementation begins. If
you disagree with an \textbf{[axiom]}, a separate alignment round is
required.}

\end{document}
